% GENERATED FILE. Edit canonical lessons, then run npm run generate:books.
% canonical-source-hash: fnv1a64:ef597b0d889f1a7a
% canonical-lessons: GU-C523-pharsan, GU-C523-khaman, GU-C523-udhiyu, GU-C523-bhavvu, GU-C523-bolvu-dip

\chapter{A Gujarati Kitchen}
\label{ch:gu-a-gujarati-kitchen}
\hlchaptermodality{523}

\begin{chapteropening}
\textbf{By the end of this chapter:} \emph{I can talk about Gujarati snacks and dishes I like.}

Complete the last lesson of chapter 523: I can talk about Gujarati snacks and dishes I like.
\end{chapteropening}

\section[\texorpdfstring{\gu{ફરસાણ} (pharsāṇ)}{ફરસાણ (pharsāṇ)}]{\gu{ફરસાણ} (pharsāṇ) — savoury snacks}

\label{lesson:GU-C523-pharsan}



\begin{quote}
Before the new one: say the Gujarati for to soak, then the Gujarati for to mill grain.
\end{quote}

\subsection*{You'll want to know: \gu{ફરસાણ}}
\textbf{\gu{ફરસાણ}} — \emph{pharsāṇ} — \textquotedblleft{}savoury snacks\textquotedblright{}.

\begin{quote}
\gu{દિવાળીમાં} \gu{અમે} \gu{ઘરે} \gu{ઘણું} \gu{ફરસાણ} \gu{બનાવીએ} \gu{છીએ}.
\end{quote}

\begin{quote}
\emph{divāḷīmā̃ ame ghare ghaṇũ pharsāṇ banāvīe chīe.}
\end{quote}

\begin{quote}
At Diwali we make lots of savoury snacks at home.
\end{quote}

A neuter noun.

\begin{cousinweb}
\textbf{\gu{ફરસાણ}} (\emph{pharsāṇ}) is the savoury side of a Gujarati meal or a festival tray. No dictionary we use records an older form, so its history is not settled.
\end{cousinweb}

\begin{grammarlens}[title={What you've built}]
One more word for describing what you lived and felt.
\end{grammarlens}

\subsection*{Your turn}
\begin{itemize}
\raggedright
  \item \emph{Say it:} \emph{pharsāṇ}
  \item \emph{Say it:} \emph{pharsāṇ}, once more
  \item \emph{Say it:} \emph{divāḷīmā̃ ame ghare ghaṇũ pharsāṇ banāvīe chīe.} — At Diwali we make lots of savoury snacks at home.
  \item \emph{Recall:} say the Gujarati for to soak, then the Gujarati for to mill grain, then say \emph{pharsāṇ} again
\end{itemize}

\begin{tcolorbox}[breakable,colback=teal!4,colframe=teal!35!black,title={Before you move on}]
What does \gu{ફરસાણ} mean? (\textquotedblleft{}Savoury snacks\textquotedblright{}.) Say it once more.
\end{tcolorbox}
\section[\texorpdfstring{\gu{ખમણ} (khamaṇ)}{ખમણ (khamaṇ)}]{\gu{ખમણ} (khamaṇ) — khaman (a steamed gram-flour snack)}

\label{lesson:GU-C523-khaman}



\begin{quote}
Before the new one: say the Gujarati for to mill grain, then the Gujarati for savoury snacks.
\end{quote}

\subsection*{You'll want to know: \gu{ખમણ}}
\textbf{\gu{ખમણ}} — \emph{khamaṇ} — \textquotedblleft{}khaman (a steamed gram-flour snack)\textquotedblright{}.

\begin{quote}
\gu{ખમણ} \gu{ચણાના} \gu{લોટમાંથી} \gu{બને} \gu{છે}.
\end{quote}

\begin{quote}
\emph{khamaṇ caṇānā loṭmā̃thī bane che.}
\end{quote}

\begin{quote}
Khaman is made from gram flour.
\end{quote}

A neuter noun.

\begin{cousinweb}
\textbf{\gu{ખમણ}} (\emph{khamaṇ}) is a soft steamed cake of gram flour, cousin to \textbf{\gu{ઢોકળાં}}. No dictionary we use records an older form, so its history is not settled.
\end{cousinweb}

\begin{grammarlens}[title={What you've built}]
One more word for describing what you lived and felt.
\end{grammarlens}

\subsection*{Your turn}
\begin{itemize}
\raggedright
  \item \emph{Say it:} \emph{khamaṇ}
  \item \emph{Say it:} \emph{khamaṇ}, once more
  \item \emph{Say it:} \emph{khamaṇ caṇānā loṭmā̃thī bane che.} — Khaman is made from gram flour.
  \item \emph{Recall:} say the Gujarati for to mill grain, then the Gujarati for savoury snacks, then say \emph{khamaṇ} again
\end{itemize}

\begin{tcolorbox}[breakable,colback=teal!4,colframe=teal!35!black,title={Before you move on}]
What does \gu{ખમણ} mean? (\textquotedblleft{}Khaman (a steamed gram-flour snack)\textquotedblright{}.) Say it once more.
\end{tcolorbox}
\section[\texorpdfstring{\gu{ઊંધિયું} (ū̃dhiyũ)}{ઊંધિયું (ū̃dhiyũ)}]{\gu{ઊંધિયું} (ū̃dhiyũ) — undhiyu (a mixed winter vegetable dish)}

\label{lesson:GU-C523-udhiyu}



\begin{quote}
Before the new one: say the Gujarati for savoury snacks, then the Gujarati for khaman.
\end{quote}

\subsection*{You'll want to know: \gu{ઊંધિયું}}
\textbf{\gu{ઊંધિયું}} — \emph{ū̃dhiyũ} — \textquotedblleft{}undhiyu (a mixed winter vegetable dish)\textquotedblright{}.

\begin{quote}
\gu{શિયાળામાં} \gu{અમારા} \gu{ઘરે} \gu{ઊંધિયું} \gu{બને} \gu{છે}.
\end{quote}

\begin{quote}
\emph{śiyāḷāmā̃ amārā ghare ū̃dhiyũ bane che.}
\end{quote}

\begin{quote}
In winter undhiyu is made in our house.
\end{quote}

A neuter noun.

\begin{cousinweb}
\textbf{\gu{ઊંધિયું}} (\emph{ū̃dhiyũ}) is built on \emph{ū̃dhũ}, \textquotedblleft{}upside down\textquotedblright{}, which CDIAL files under Sanskrit \emph{avamūrdha}, \textquotedblleft{}head downwards\textquotedblright{} (CDIAL 804); the dish is said to take its name from the pot once cooked upside down in the ground.
\end{cousinweb}

\begin{grammarlens}[title={What you've built}]
One more word for describing what you lived and felt.
\end{grammarlens}

\subsection*{Your turn}
\begin{itemize}
\raggedright
  \item \emph{Say it:} \emph{ū̃dhiyũ}
  \item \emph{Say it:} \emph{ū̃dhiyũ}, once more
  \item \emph{Say it:} \emph{śiyāḷāmā̃ amārā ghare ū̃dhiyũ bane che.} — In winter undhiyu is made in our house.
  \item \emph{Recall:} say the Gujarati for savoury snacks, then the Gujarati for khaman, then say \emph{ū̃dhiyũ} again
\end{itemize}

\begin{tcolorbox}[breakable,colback=teal!4,colframe=teal!35!black,title={Before you move on}]
What does \gu{ઊંધિયું} mean? (\textquotedblleft{}Undhiyu (a mixed winter vegetable dish)\textquotedblright{}.) Say it once more.
\end{tcolorbox}
\section[\texorpdfstring{\gu{ભાવવું} (bhāvvũ)}{ભાવવું (bhāvvũ)}]{\gu{ભાવવું} (bhāvvũ) — to be to one's taste}

\label{lesson:GU-C523-bhavvu}



\begin{quote}
Before the new one: say the Gujarati for khaman, then the Gujarati for undhiyu.
\end{quote}

\subsection*{You'll want to know: \gu{ભાવવું}}
\textbf{\gu{ભાવવું}} — \emph{bhāvvũ} — \textquotedblleft{}to be to one's taste\textquotedblright{}.

\begin{quote}
\gu{મને} \gu{ઊંધિયું} \gu{ઘણું} \gu{ભાવે} \gu{છે}.
\end{quote}

\begin{quote}
\emph{mane ū̃dhiyũ ghaṇũ bhāve che.}
\end{quote}

\begin{quote}
I really love undhiyu.
\end{quote}

Past (it …): \gu{તે} \gu{શાક} \gu{તેને} \gu{ભાવ્યું}, \emph{te śāk tene bhāvyũ}.

\begin{cousinweb}
\textbf{\gu{ભાવવું}} (\emph{bhāvvũ}) is filed by CDIAL under Sanskrit \emph{bhāti}, \textquotedblleft{}shines\textquotedblright{}, and CDIAL lists the Gujarati verb, \textquotedblleft{}to please, be agreeable to\textquotedblright{} (CDIAL 9445). Said of food, the eater takes \textbf{\gu{ને}} and the dish is the subject: \textbf{\gu{મને} \gu{ઊંધિયું} \gu{ભાવે} \gu{છે}}.
\end{cousinweb}

\begin{grammarlens}[title={What you've built}]
One more word for describing what you lived and felt.
\end{grammarlens}

\subsection*{Your turn}
\begin{itemize}
\raggedright
  \item \emph{Say it:} \emph{bhāvvũ}
  \item \emph{Say it:} \emph{bhāvvũ}, once more
  \item \emph{Say it:} \emph{mane ū̃dhiyũ ghaṇũ bhāve che.} — I really love undhiyu.
  \item \emph{Recall:} say the Gujarati for khaman, then the Gujarati for undhiyu, then say \emph{bhāvvũ} again
\end{itemize}

\begin{tcolorbox}[breakable,colback=teal!4,colframe=teal!35!black,title={Before you move on}]
What does \gu{ભાવવું} mean? (\textquotedblleft{}To be to one's taste\textquotedblright{}.) Say it once more.
\end{tcolorbox}
\section[\texorpdfstring{\gu{બોળવું} (boḷvũ)}{બોળવું (boḷvũ)}]{\gu{બોળવું} (boḷvũ) — to dunk (in a liquid)}

\label{lesson:GU-C523-bolvu-dip}



\begin{quote}
Before the new one: say the Gujarati for undhiyu, then the Gujarati for to be to one's taste.
\end{quote}

\subsection*{You'll want to know: \gu{બોળવું}}
\textbf{\gu{બોળવું}} — \emph{boḷvũ} — \textquotedblleft{}to dunk (in a liquid)\textquotedblright{}.

\begin{quote}
\gu{મીરાએ} \gu{ભજિયાં} \gu{ચટણીમાં} \gu{બોળ્યાં}.
\end{quote}

\begin{quote}
\emph{Mīrāe bhajiyā̃ caṭṇīmā̃ boḷyā̃.}
\end{quote}

\begin{quote}
Mira dipped the fritters in the chutney.
\end{quote}

Past, with \gu{તેણે} (he or she …): \gu{તેણે} \gu{રોટલી} \gu{બોળી}, \emph{teṇe roṭlī boḷī}.

\begin{cousinweb}
\textbf{\gu{બોળવું}} (\emph{boḷvũ}) comes from \emph{*buḍyati}, \textquotedblleft{}sinks\textquotedblright{}, a form CDIAL reconstructs, and CDIAL lists the Gujarati verb, \textquotedblleft{}to dip\textquotedblright{} (CDIAL 9272).
\end{cousinweb}

\begin{grammarlens}[title={What you've built}]
One more word for describing what you lived and felt.
\end{grammarlens}

\subsection*{Your turn}
\begin{itemize}
\raggedright
  \item \emph{Say it:} \emph{boḷvũ}
  \item \emph{Say it:} \emph{boḷvũ}, once more
  \item \emph{Say it:} \emph{Mīrāe bhajiyā̃ caṭṇīmā̃ boḷyā̃.} — Mira dipped the fritters in the chutney.
  \item \emph{Recall:} say the Gujarati for undhiyu, then the Gujarati for to be to one's taste, then say \emph{boḷvũ} again
\end{itemize}

\begin{tcolorbox}[breakable,colback=teal!4,colframe=teal!35!black,title={Before you move on}]
What does \gu{બોળવું} mean? (\textquotedblleft{}To dunk (in a liquid)\textquotedblright{}.) Say it once more.
\end{tcolorbox}
